\section{Study Area}
\label{sec:study_area}

\subsection{Geographic Setting}
\label{sec:geography}

Himachal Pradesh occupies an area of approximately 55,673\,km$^2$
in the north-western Indian Himalaya, extending from
approximately 30.22°N to 33.17°N latitude and 75.47°E to 79.04°E longitude
(\cref{fig:study_area}).
The state spans an extreme elevational range, from approximately 350\,m\,asl
in the Shivalik foothills of Sirmaur district to 6,816\,m\,asl
at the peak of Reo Purgyil in Kinnaur, producing extraordinary climatic diversity
within a relatively compact geographical extent.

Five major river systems drain HP:
(1) the \textbf{Beas}, draining the Kullu valley and Rohtang plateau (3,978\,km$^2$
contribution from HP);
(2) the \textbf{Satluj} (Sutlej), the largest river system, draining the Spiti,
Kinnaur, and Shimla regions;
(3) the \textbf{Chenab}, draining Lahaul and upper Chamba (receiving substantial
glacier melt from the highest elevations);
(4) the \textbf{Ravi}, draining central and lower Chamba; and
(5) the \textbf{Yamuna}, draining the south-eastern districts of Sirmaur and Solan.
All five systems ultimately join the Indus basin.
These five basins form the spatial blocks used in cross-validation
(Section~\ref{sec:validation}).

\subsection{Physiographic Zones}
\label{sec:physiography}

HP is conventionally divided into three physiographic zones
that exhibit distinct flash flood regimes:

\textbf{Shivalik Zone (Sub-Himalayan, 350--1500\,m):}
The south-facing, densely forested Shivalik ranges experience
high-intensity monsoonal rainfall (1,200--2,500\,mm/year).
Flash floods here are predominantly rainfall-triggered,
with rapid runoff generation on saturated soils during
late-July to mid-August cloudburst events.
Districts: Bilaspur, Una, Hamirpur, lower Kangra.

\textbf{Lesser and Greater Himalaya (1500--4500\,m):}
The temperate mid-Himalayan belt encompasses Kullu, Mandi, Chamba, Sirmaur, and Shimla.
This zone has the most complex flash flood triggering regime:
monsoonal rainfall (900--2,000\,mm/year), antecedent snowmelt,
and landslide-dam outburst events all contribute.
The Beas and upper Satluj valleys carry the highest documented flash flood density
in HP, as evidenced by the HiFlo-DAT database of 128 events in Kullu alone
\citep{hiflodot2025}.

\textbf{Trans-Himalayan Zone (Lahaul-Spiti, Kinnaur, upper Chamba, 2000--6800\,m):}
A cold, high-altitude semi-arid zone receiving less than 500\,mm annual precipitation,
mostly as snowfall.
Flash floods in this zone are predominantly triggered by Glacial Lake Outburst Floods (GLOFs),
snowmelt pulses, and rare but intense convective events.
The rapid expansion of glacial lakes (area +75\% in Chenab basin 1990--2022;
\citealt{chenab2024glof}) is producing a new and accelerating GLOF risk
that existing susceptibility maps do not capture.

\subsection{Hydroclimatology}
\label{sec:climate}

HP's climate is controlled by three competing moisture sources:
(1) the Indian Summer Monsoon (ISM, June--September), which dominates in the
southern and central districts;
(2) western disturbances (December--March), bringing snowfall to higher elevations; and
(3) local convection, particularly in the form of cloudbursts,
which are intense, localised rainfall events ($>$100\,mm\,h$^{-1}$)
concentrated in July--August.
The compound interaction of monsoon rainfall, snowmelt-induced elevated soil moisture,
and cloudburst events was identified as the primary driver of the catastrophic
2023 flood season by \citet{dixit2026nhess} and is explicitly captured by
the antecedent precipitation conditioning factor included in this study.

Mean annual rainfall ranges from approximately 450\,mm in the Spiti valley
to over 3,200\,mm in parts of Dharamshala (Kangra) and the Shivalik fringe,
making HP one of the most rainfall-diverse states in India.
Extreme rainfall quantiles (95th percentile of daily event rainfall)
are spatially heterogeneous and are included as a separate conditioning factor
to capture the episodic, high-intensity triggering mechanism of flash floods
rather than mean annual conditions alone.

\subsection{Flood History and Significance}
\label{sec:flood_history}

Flash floods and cloudbursts have caused systematic losses in HP throughout
the recorded period.
The HiFlo-DAT database \citep{hiflodot2025}, covering Kullu district from 1900 to 2023,
documents 128 significant flood events.
Major recent events include:
the 2018 Beas flash flood (Kullu--Manali highway, 8 deaths);
the 2021 Spiti GLOF (Sumdo, 11 deaths);
the 2022 Beas cloudburst series (Kullu, 22 deaths);
and the 2023 multi-basin monsoon floods (HP-wide, 404 deaths, Rs.\,9,905\,crore;
\citealt{ndma2023}).
The 2023 events are used as an independent temporal test set in this study
(Section~\ref{sec:validation}).

HP supports approximately 17.17 million residents and receives
6--7 million tourists annually, with peak tourist concentration
in July--August coinciding precisely with the peak flash flood season.
The Manali--Leh National Highway (NH-3) and the Chandigarh--Manali National Highway (NH-21)
are critical transport arteries that are repeatedly cut by flash floods;
HP hosts 27 major hydroelectric projects with a combined installed capacity of
10,569\,MW \citep{niti2022hp}, most located in high-susceptibility valley corridors.
Mapping and communicating flash flood risk in these corridors is therefore
an urgent public safety priority.
